\section{Diagramas de Estado}\label{sec:estado}

Nas figuras a seguir são apresentados os diagramas de estado das entidades do sistema,
descrevendo os ciclos de vida e as transições de estado possíveis.

\subsection{Usuário}

Na Figura~\ref{fig:estado-user} é representado o ciclo de vida da entidade \texttt{User} do módulo de gerenciamento de usuários.

\begin{figure}[H]
    \caption{Diagrama de estado do usuário}
    \label{fig:estado-user}
    \centering
    \includegraphics[width=0.95\textwidth,height=0.75\textheight,keepaspectratio]{../../../out/src/plantuml/estado/user_status.png}

    {\footnotesize \textbf{Fonte:} O autor.}
\end{figure}

\subsection{Sala}

Na Figura~\ref{fig:estado-room}, Figura~\ref{fig:estado-room-type}, Figura~\ref{fig:estado-building} e Figura~\ref{fig:estado-resource} são representados respectivamente os ciclos de vida das entidades \texttt{Room}, \texttt{RoomType}, \texttt{Building} e \texttt{Resource} do módulo de gerenciamento de salas.

\begin{figure}[H]
    \caption{Diagrama de estado da sala}
    \label{fig:estado-room}
    \centering
    \includegraphics[width=0.95\textwidth,height=0.75\textheight,keepaspectratio]{../../../out/src/plantuml/estado/room_status.png}

    {\footnotesize \textbf{Fonte:} O autor.}
\end{figure}

\begin{figure}[H]
    \caption{Diagrama de estado do tipo de sala}
    \label{fig:estado-room-type}
    \centering
    \includegraphics[width=0.95\textwidth,height=0.75\textheight,keepaspectratio]{../../../out/src/plantuml/estado/room_type_status.png}

    {\footnotesize \textbf{Fonte:} O autor.}
\end{figure}

\begin{figure}[H]
    \caption{Diagrama de estado do prédio}
    \label{fig:estado-building}
    \centering
    \includegraphics[width=0.95\textwidth,height=0.75\textheight,keepaspectratio]{../../../out/src/plantuml/estado/building_status.png}

    {\footnotesize \textbf{Fonte:} O autor.}
\end{figure}

\begin{figure}[H]
    \caption{Diagrama de estado do recurso}
    \label{fig:estado-resource}
    \centering
    \includegraphics[width=0.95\textwidth,height=0.75\textheight,keepaspectratio]{../../../out/src/plantuml/estado/resource_status.png}

    {\footnotesize \textbf{Fonte:} O autor.}
\end{figure}

\subsection{Evento}

Na Figura~\ref{fig:estado-event}, Figura~\ref{fig:estado-event-request} e Figura~\ref{fig:estado-period} são representados respectivamente os ciclos de vida das entidades \texttt{Event}, \texttt{EventRequest} e \texttt{Period} do módulo de gerenciamento de eventos.

\begin{figure}[H]
    \caption{Diagrama de estado do evento}
    \label{fig:estado-event}
    \centering
    \includegraphics[width=0.95\textwidth,height=0.75\textheight,keepaspectratio]{../../../out/src/plantuml/estado/event_status.png}

    {\footnotesize \textbf{Fonte:} O autor.}
\end{figure}

\begin{figure}[H]
    \caption{Diagrama de estado da solicitação de alteração}
    \label{fig:estado-event-request}
    \centering
    \includegraphics[width=0.95\textwidth,height=0.75\textheight,keepaspectratio]{../../../out/src/plantuml/estado/event_request_status.png}

    {\footnotesize \textbf{Fonte:} O autor.}
\end{figure}

\begin{figure}[H]
    \caption{Diagrama de estado do período}
    \label{fig:estado-period}
    \centering
    \includegraphics[width=0.95\textwidth,height=0.75\textheight,keepaspectratio]{../../../out/src/plantuml/estado/period_status.png}

    {\footnotesize \textbf{Fonte:} O autor.}
\end{figure}

\subsection{Relatório}

Na Figura~\ref{fig:estado-report} é representado o ciclo de vida da entidade \texttt{Report} do módulo de relatórios.

\begin{figure}[H]
    \caption{Diagrama de estado do relatório}
    \label{fig:estado-report}
    \centering
    \includegraphics[width=0.95\textwidth,height=0.75\textheight,keepaspectratio]{../../../out/src/plantuml/estado/report_status.png}

    {\footnotesize \textbf{Fonte:} O autor.}
\end{figure}
